\documentclass[aspectratio=169]{beamer}
\input{header}
\coursecode{JKSSB}

\title{Current Multimedia, GUI \& Media Formats}
\subtitle{Lesson 65 --- Multimedia, GUI and File-Format Recall}
\author{JKSSB Computer Awareness}
\date{\today}

\begin{document}
\frame{\titlepage}

\begin{frame}{What Is Multimedia?}
  \begin{itemize}
    \item \key{Multimedia} is content that combines more than one medium in a
      single presentation.
    \item The media may be \key{text, graphics, audio, video}, and
      \key{animation}.
    \item Multimedia is stored in digital form, so each medium must first be
      captured and converted into binary data.
    \item Examples: a training video with narration, a web page with pictures,
      and a presentation with sound.
  \end{itemize}
  \begin{block}{The one idea}
    Multimedia = multiple media types presented together, all reduced to
    digital form the computer can process.
  \end{block}
\end{frame}

\begin{frame}{The Five Building Blocks of Multimedia}
  \begin{center}
    \begin{tabular}{p{0.20\textwidth}p{0.68\textwidth}}
      \toprule
      \textbf{Element} & \textbf{Meaning} \\
      \midrule
      Text & Readable characters; the base medium. \\
      Graphics & Still images, drawings, and photographs. \\
      Audio & Sound, music, and speech. \\
      Video & Moving pictures, usually with sound. \\
      Animation & A sequence of images that creates motion. \\
      \bottomrule
    \end{tabular}
  \end{center}
  \vspace{0.25cm}
  \muted{Text and graphics are static; audio, video, and animation are
  time-based media.}
\end{frame}

\begin{frame}{Digitising Media: Sampling}
  \begin{itemize}
    \item Analogue sound and light are \key{continuous} signals; a computer
      stores only discrete numbers.
    \item \key{Sampling} means measuring the signal at regular intervals and
      recording each measurement.
    \item The \key{sampling rate} is the number of samples taken per second,
      measured in hertz (Hz).
    \item A \key{higher} sampling rate gives better quality but a larger file.
  \end{itemize}
  \begin{block}{Sampling answers two questions}
    How often do we measure the signal, and how many measurements do we store
    each second?
  \end{block}
\end{frame}

\begin{frame}{Quantization and Bit Depth}
  \begin{itemize}
    \item Each sample is a precise value; the computer must round it to the
      nearest allowed level.
    \item \key{Quantization} is the process of mapping each sample to one of a
      fixed set of discrete levels.
    \item The number of levels is set by the \key{bit depth} of each sample.
    \item \key{More levels} (higher bit depth) means less rounding error and
      better quality, but a larger file.
  \end{itemize}
  \begin{alertblock}{Trap alert}
    Sampling = how often we measure. Quantization = how precisely each
    measurement is recorded. Do not confuse the two.
  \end{alertblock}
\end{frame}

\begin{frame}{Compression: Lossy vs Lossless}
  \begin{itemize}
    \item \key{Compression} reduces the size of a media file.
    \item \key{Lossless} compression restores the original exactly --- no data
      is lost (example: PNG, GIF, ZIP).
    \item \key{Lossy} compression discards some data to shrink the file further
      (examples: JPEG, MP3).
    \item Lossy files are usually smaller; lossless files preserve full quality.
  \end{itemize}
  \begin{block}{Which to choose}
    Use lossless when exactness matters; use lossy when a smaller file matters
    more than perfect fidelity.
  \end{block}
  \begin{alertblock}{Exam point}
    JPEG, MP3, and MP4 are lossy formats; PNG and GIF are lossless.
  \end{alertblock}
\end{frame}

\begin{frame}{Synchronization of Media Streams}
  \begin{itemize}
    \item In a multimedia file, audio and video are separate \key{streams}.
    \item \key{Synchronization} keeps those streams aligned in time.
    \item If they drift apart, speech no longer matches the moving lips ---
      a \key{sync} error.
    \item Synchronization is what makes a video play with the sound in the
      correct place.
  \end{itemize}
  \begin{block}{Simple view}
    Synchronization means the picture and the sound stay together at every
    moment of playback.
  \end{block}
\end{frame}

\begin{frame}{GUI vs CLI}
  \begin{itemize}
    \item A \key{GUI} (Graphical User Interface) lets the user work with
      \key{pictures, icons, and menus} instead of typed commands.
    \item A \key{CLI} (Command-Line Interface) accepts typed text commands.
    \item The GUI is easier for beginners; the CLI is often faster for experts.
    \item Windows and modern Linux desktops use a GUI.
  \end{itemize}
  \begin{alertblock}{Trap alert}
    GUI = \key{Graphical} User Interface. It is not ``General'' or ``Graphic''
    alone; the exact full form is asked directly.
  \end{alertblock}
\end{frame}

\begin{frame}{GUI Elements: WIMP}
  \begin{itemize}
    \item The four classic GUI elements are remembered as \key{WIMP}:
      \key{W}indows, \key{I}cons, \key{M}enus, and \key{P}ointers.
    \item \key{Window}: a rectangular area that holds a program or document.
    \item \key{Icon}: a small picture that represents a file, folder, or
      program.
    \item \key{Menu}: a list of commands offered by the software.
    \item \key{Pointer}: the on-screen arrow moved by the mouse or touchpad.
  \end{itemize}
  \begin{block}{WIMP full form}
    Windows, Icons, Menus, Pointers --- the standard interaction set of a GUI.
  \end{block}
\end{frame}

\begin{frame}{Dialogue Boxes and WYSIWYG}
  \begin{itemize}
    \item A \key{dialogue box} is a small window that asks for input or
      confirms an action (for example, Save As or Print).
    \item It usually contains text boxes, drop-down lists, check boxes, radio
      buttons, and \key{OK / Cancel} buttons.
    \item \key{WYSIWYG} (What You See Is What You Get) means the screen shows
      the document exactly as it will be printed.
  \end{itemize}
  \begin{exampleblock}{Worked example}
    In MS Word, the page on screen looks like the printed page. That is
    WYSIWYG. When the Print dialogue box opens, the settings appear inside a
    dialogue box.
  \end{exampleblock}
\end{frame}

\begin{frame}{Image Formats: PNG, JPEG, GIF}
  \begin{center}
    \begin{tabular}{p{0.13\textwidth}p{0.24\textwidth}p{0.50\textwidth}}
      \toprule
      \textbf{Format} & \textbf{Full form} & \textbf{Nature and typical use} \\
      \midrule
      PNG & Portable Network Graphics & Lossless; supports transparency; logos, web graphics. \\
      JPEG & Joint Photographic Experts Group & Lossy; best for photographs. \\
      GIF & Graphics Interchange Format & Lossless but limited to 256 colours; simple animation. \\
      \bottomrule
    \end{tabular}
  \end{center}
  \vspace{0.2cm}
  \muted{GIF is an image format that also supports short, simple animation.}
\end{frame}

\begin{frame}{Choosing the Right Image Format}
  \begin{itemize}
    \item \key{JPEG} for photographs and complex pictures --- small size,
      some quality loss.
    \item \key{PNG} when transparency or exact colours are needed --- lossless.
    \item \key{GIF} for simple graphics and small animations --- limited to
      256 colours.
    \item \key{BMP} is a large, uncompressed bitmap format, not a web format.
  \end{itemize}
  \begin{alertblock}{Trap alert}
    JPEG is \key{lossy}, PNG and GIF are \key{lossless}. Neither PNG nor JPEG
    supports animation; only GIF does among the three.
  \end{alertblock}
\end{frame}

\begin{frame}{Audio Format: MP3}
  \begin{itemize}
    \item \key{MP3} stands for \key{MPEG-1 Audio Layer III} (also written
      MPEG Audio Layer III).
    \item It is a \key{lossy} audio format: it removes sound data the ear is
      unlikely to notice.
    \item Its small size made it the standard format for portable music and
      downloads.
    \item \key{WAV} is an uncompressed audio format; MP3 is compressed.
  \end{itemize}
  \begin{block}{Keep the family straight}
    MPEG is the Moving Picture Experts Group; MP3 is its audio layer, MP4 is
    its later container.
  \end{block}
\end{frame}

\begin{frame}{Video Format: MP4}
  \begin{itemize}
    \item \key{MP4} stands for \key{MPEG-4 Part 14}.
    \item It is a \key{container} format: it can hold video, audio, and
      subtitles in one file.
    \item It is widely used for streaming, downloads, and mobile video.
    \item \key{AVI} (Audio Video Interleave) is an older video container.
  \end{itemize}
  \begin{block}{Container vs codec}
    MP4 is the container (the box); the codec inside it actually compresses
    the video and audio.
  \end{block}
\end{frame}

\begin{frame}{Media Format Quick Table}
  \begin{center}
    \begin{tabular}{p{0.12\textwidth}p{0.20\textwidth}p{0.18\textwidth}p{0.36\textwidth}}
      \toprule
      \textbf{Format} & \textbf{Type} & \textbf{Compression} & \textbf{Best use} \\
      \midrule
      PNG & Image & Lossless & Web graphics, transparency \\
      JPEG & Image & Lossy & Photographs \\
      GIF & Image & Lossless & Simple graphics, animation \\
      MP3 & Audio & Lossy & Music and speech \\
      MP4 & Video & Lossy & Video with sound \\
      \bottomrule
    \end{tabular}
  \end{center}
  \vspace{0.2cm}
  \muted{Full forms are high-yield: learn every expansion exactly.}
\end{frame}

\begin{frame}{Plotter}
  \begin{itemize}
    \item A \key{plotter} is an \key{output} device that draws on paper.
    \item It produces \key{vector graphics} --- lines and curves defined by
      coordinates.
    \item It is used for \key{large engineering drawings}, maps, blueprints,
      and banners.
    \item Pen plotters use pens; modern inkjet plotters spray ink.
  \end{itemize}
  \begin{alertblock}{Trap alert}
    A plotter is an \key{output} device, not input. It is chosen for large
    technical or engineering drawings, not for ordinary text. (PYQ-29)
  \end{alertblock}
\end{frame}

\begin{frame}{Impact vs Non-Impact Printers (Revisit)}
  \begin{center}
    \begin{tabular}{p{0.16\textwidth}p{0.36\textwidth}p{0.36\textwidth}}
      \toprule
      \textbf{Point} & \textbf{Impact} & \textbf{Non-impact} \\
      \midrule
      Action & Strikes the paper & No striking \\
      Examples & Dot-matrix, daisy-wheel, line printer & Inkjet, laser, thermal \\
      Copies & Carbon copies possible & No carbon copies \\
      Noise & Noisy & Quiet \\
      \bottomrule
    \end{tabular}
  \end{center}
  \vspace{0.2cm}
  \muted{Both are output devices; the difference is whether the mechanism
  strikes the paper.}
\end{frame}

\begin{frame}{Full Forms You Must Recall}
  \begin{center}
    \begin{tabular}{p{0.16\textwidth}p{0.40\textwidth}p{0.30\textwidth}}
      \toprule
      \textbf{Short} & \textbf{Full form} & \textbf{Note} \\
      \midrule
      GUI & Graphical User Interface & Icons and menus, not typed commands \\
      CLI & Command-Line Interface & Typed commands \\
      WYSIWYG & What You See Is What You Get & Screen matches print \\
      WIMP & Windows, Icons, Menus, Pointers & Classic GUI elements \\
      MP3 & MPEG-1 Audio Layer III & Lossy audio \\
      MP4 & MPEG-4 Part 14 & Video container \\
      \bottomrule
    \end{tabular}
  \end{center}
\end{frame}

\begin{frame}{Trap Alert: Format Facts to Keep Straight}
  \begin{itemize}
    \item \key{Sampling} = how often we measure; \key{quantization} = how
      precisely each sample is stored.
    \item \key{JPEG, MP3, MP4} are lossy; \key{PNG, GIF} are lossless.
    \item \key{MP3} is audio; \key{MP4} is video; do not swap their full forms.
    \item A \key{plotter} is an output device for large drawings.
    \item \key{Impact} printers strike the paper; \key{non-impact} printers do
      not.
    \item A \key{dialogue box} asks for input; \key{WYSIWYG} shows the print
      layout on screen.
  \end{itemize}
\end{frame}

\begin{frame}{Lesson 65: Recall}
  \begin{itemize}
    \item Multimedia combines text, graphics, audio, video, and animation.
    \item \key{Sampling} measures a continuous signal; \key{quantization} maps
      each sample to discrete levels; \key{bit depth} sets the number of
      levels.
    \item \key{Lossless} restores the original (PNG, GIF); \key{lossy} discards
      data (JPEG, MP3).
    \item \key{Synchronization} keeps audio and video lined up in time.
    \item GUI elements are \key{Windows, Icons, Menus, Pointers} (WIMP);
      \key{WYSIWYG} matches screen to print.
    \item \key{PNG}, \key{JPEG}, \key{GIF} are image formats; \key{MP3} is
      audio; \key{MP4} is video.
    \item A \key{plotter} is an output device for large drawings; \key{impact}
      printers strike, \key{non-impact} printers do not.
  \end{itemize}
\end{frame}
\end{document}
